% ============================================================
% TWO PROPOSED METHODS AND CONTROLLED ABLATION
% ============================================================

\paragraph{Proposed Plain Fourier U-Net.}
The first proposed method, denoted by \texttt{plain\_fourier\_unet},
introduces an identity-initialized Fourier bottleneck into the U-Net
architecture. The bottleneck jointly learns bounded amplitude and phase
corrections while retaining a residual spatial path.

\paragraph{Proposed APDR-Fourier U-Net.}
The second proposed method, denoted by \texttt{apdr\_fourier\_unet}, extends
the complete Plain Fourier U-Net path by a zero-initialized
amplitude--phase disagreement-routed residual adapter. Thus,
\[
\mathrm{APDR\text{-}Fourier\ U\text{-}Net}
=
\mathrm{Plain\ Fourier\ U\text{-}Net}
+
\mathrm{APDR\ residual\ adapter}.
\]
Because the residual gate is initialized at zero, both proposed methods produce
identical logits at initialization. The controlled ablation uses the same data
split, augmentation, optimizer, learning rate, loss, number of epochs,
threshold, and random seeds; hence, it isolates the contribution of the APDR
adapter.

\begin{table}[htbp]
\centering
\caption{Controlled ablation of the two proposed Fourier U-Net methods.}
\label{tab:two-proposal-ablation}
\begin{tabular}{lcc}
\hline
Method & Plain Fourier path & APDR residual adapter \\
\hline
Plain Fourier U-Net (Proposal I) & \checkmark & -- \\
APDR-Fourier U-Net (Proposal II) & \checkmark & \checkmark \\
\hline
\end{tabular}
\end{table}

% ============================================================
% RETAINED CHALLENGE DATASET
% ============================================================

\paragraph{ISBI-2012 challenge.}
The ISBI-2012 neuronal-structure electron-microscopy dataset is retained as a
supported challenge dataset. Since the public challenge test volume does not
provide ground-truth labels, the supervised experiments use the 30 labeled
training slices and construct a reproducible contiguous train--validation--test
split to reduce leakage between adjacent slices.
